% © 2026 Ing. David Jaroš — CC BY-NC-ND 4.0
%
% This work is licensed under a Creative Commons Attribution-NonCommercial-NoDerivatives
% 4.0 International License (CC BY-NC-ND 4.0).
%
% License History: Earlier drafts (up to v0.3) were released under CC BY 4.0.
% From v0.4 onward, all material is released under CC BY-NC-ND 4.0 to protect
% the integrity of the theoretical work during ongoing academic development.
%
% See LICENSE.md for full license text.

\ifdefined\INCLUDEMODE
\else
  \documentclass[12pt,a4paper]{article}
  \usepackage[a4paper, margin=2.5cm]{geometry}
  \usepackage{amsmath, amssymb, amsthm}
  \usepackage{mathtools}
  \usepackage{hyperref}
  \usepackage{xcolor}
  \usepackage{mdframed}
  \usepackage{booktabs}

  \newtheorem{theorem}{Theorem}[section]
  \newtheorem{lemma}[theorem]{Lemma}
  \newtheorem{proposition}[theorem]{Proposition}
  \newtheorem{conjecture}[theorem]{Conjecture}
  \theoremstyle{definition}
  \newtheorem{definition}[theorem]{Definition}
  \theoremstyle{remark}
  \newtheorem{remark}[theorem]{Remark}

  \newmdenv[backgroundcolor=yellow!20, linecolor=orange, linewidth=1.5pt]{gapbox}
  \newmdenv[backgroundcolor=green!15, linecolor=green!60!black, linewidth=1.5pt]{resultbox}
  \newmdenv[backgroundcolor=red!10, linecolor=red!60, linewidth=1.5pt]{deadendbox}
  \newmdenv[backgroundcolor=blue!8, linecolor=blue!50, linewidth=1.5pt]{insightbox}

  \title{\textbf{Gap A --- Approach A2: Prime Geometry and Modular Curves}\\
         \large Can $B_{\mathrm{phenom}} \approx 46.3$ Arise from
         the Geometry of $X_0(137)$ or $X_0(139)$?}
  \author{David Jaro\v{s}}
  \date{2026-03-06}

  \begin{document}
  \maketitle

  \begin{abstract}
  We investigate whether the phenomenological coefficient
  $B_{\mathrm{phenom}} \approx 46.3$ is determined by the geometry of the
  modular curves $X_0(137)$ and $X_0(139)$.  We compute the genera, volumes,
  and natural period lengths of these curves, and test whether any combination
  reproduces $B_{\mathrm{phenom}}$.  We find that $N_{\mathrm{eff}}^{3/2} = 12^{3/2}
  \approx 41.6$ is close but fails.  The volume of $X_0(139)$ divided by a
  natural UBT normalisation is identified as a remaining candidate.
  Results are numerically reproducible via \texttt{tools/compute\_modular\_curve\_genus.py}.
  \end{abstract}
\fi

\section{Context}
\label{sec:context}

In UBT the prime attractor theorem (proved) selects $n^* \in \{137, 139\}$
as stable winding numbers.  The modular curve $X_0(N)$ is the natural
geometric arena when the prime $p = N$ plays a distinguished role.  Its
genus $g(N)$ encodes arithmetic information about level-$N$ modular forms.

\begin{insightbox}
\textbf{Key observation:} The genus of $X_0(p)$ for primes near the UBT
attractors:
\begin{center}
\begin{tabular}{cc}
  $g(X_0(137)) = 11$, & $g(X_0(139)) = 11$.
\end{tabular}
\end{center}
Note: $N_{\mathrm{eff}} = 12$, but the genus $g(139) = 11 \neq N_{\mathrm{eff}}$.
The connection to $N_{\mathrm{eff}}$ is \emph{not} through the genus directly.
Instead, the index $\mu(\Gamma_0(137)) = 138 = 3 \times 46$ is the key quantity.
\end{insightbox}

%──────────────────────────────────────────────────────────────────────────────
\section{Genus of $X_0(N)$ for Relevant Primes}
\label{sec:genera}
%──────────────────────────────────────────────────────────────────────────────

For a prime $p$, the genus of $X_0(p)$ is given by the formula
\begin{equation}
  g(p) \;=\; \frac{p-13}{12} + \epsilon(p),
  \label{eq:genus_prime}
\end{equation}
where $\epsilon(p) \in \{0, 1/2\}$ is a correction from elliptic and
cuspidal contributions.  More precisely:
\begin{equation}
  g(p) \;=\; \left\lfloor\frac{p+1}{12}\right\rfloor
             - \left\lfloor\frac{p}{4}\right\rfloor + \left\lfloor\frac{p}{4}\right\rfloor
             + \cdots,
\end{equation}
but for our purposes the direct values suffice:

\begin{center}
\begin{tabular}{cccc}
  \toprule
  Prime $p$ & $g(X_0(p))$ & Notes \\
  \midrule
  127 & 10 & \\
  131 & 11 & \\
  137 & \textbf{11} & prime attractor \\
  139 & \textbf{11} & prime attractor \\
  149 & 12 & $= N_{\mathrm{eff}}$ (but not a prime attractor) \\
  \bottomrule
\end{tabular}
\end{center}

The full computation is in \texttt{tools/compute\_modular\_curve\_genus.py}.

%──────────────────────────────────────────────────────────────────────────────
\section{Tests Using the Genus}
\label{sec:genus_tests}
%──────────────────────────────────────────────────────────────────────────────

\subsection{Test 1: $B_{\mathrm{phenom}} = N_{\mathrm{eff}}^{3/2}$}

\begin{equation}
  N_{\mathrm{eff}}^{3/2} \;=\; 12^{3/2} \;=\; 12\sqrt{12}
    \;=\; 41.569\ldots
\end{equation}
This is close to $B_{\mathrm{phenom}} \approx 46.3$ but the relative
error is $\approx 10\%$, too large to claim an exact relation.

\begin{deadendbox}
\textbf{Dead end:} $N_{\mathrm{eff}}^{3/2} = 41.57 \neq 46.3$.
Discrepancy $\approx 10\%$.  Not a coincidence-level match.
\end{deadendbox}

\subsection{Test 2: $B_{\mathrm{phenom}} = g(139) \times \pi$}

\begin{equation}
  g(139) \times \pi \;=\; 12\pi \;\approx\; 37.70.
\end{equation}
Too small.

\subsection{Test 3: $B_{\mathrm{phenom}} = g(139) \times (1 + 1/\sqrt{g(139)}) \times \text{something}$}

No natural normalisation in UBT gives $12 \times (\text{factor}) = 46.3$
with a factor that is algebraically determined.  Ruled out by lack of
motivation.

\subsection{Test 4: Volume of $X_0(N)$}
\label{subsec:volume}

The (hyperbolic) area of the fundamental domain of $\Gamma_0(N)$ is
\begin{equation}
  \mathrm{vol}\bigl(X_0(N)\bigr)
  \;=\; \frac{\pi}{3}\,\mu(\Gamma_0(N)),
  \quad
  \mu(\Gamma_0(N)) \;=\; N\prod_{p\mid N}\!\Bigl(1+\tfrac{1}{p}\Bigr),
\end{equation}
where the product is over primes dividing $N$.  For prime $N = p$:
\begin{equation}
  \mu(\Gamma_0(p)) \;=\; p+1, \qquad
  \mathrm{vol}(X_0(p)) \;=\; \frac{\pi(p+1)}{3}.
\end{equation}

\begin{itemize}
  \item $\mathrm{vol}(X_0(137)) = \pi \times 138/3 = 46\pi \approx 144.5$.
  \item $\mathrm{vol}(X_0(139)) = \pi \times 140/3 \approx 146.6$.
\end{itemize}

Both are $\gg B_{\mathrm{phenom}} \approx 46.3$.

\begin{insightbox}
\textbf{Striking near-coincidence:}
\[
  \frac{\mathrm{vol}(X_0(137))}{\pi^2}
  \;=\; \frac{46\pi}{\pi^2}
  \;=\; \frac{46}{\pi}
  \;\approx\; 14.64.
\]
Also:
\[
  \frac{\mu(\Gamma_0(137))}{B_0}
  \;=\; \frac{138}{8\pi}
  \;\approx\; 5.49.
\]
Neither matches $B_{\mathrm{phenom}}$ directly.

However:
\[
  \frac{\mathrm{vol}(X_0(137))}{\pi \times (p-1)/6}
  \;=\; \frac{46\pi}{\pi \times 136/6}
  \;=\; \frac{46 \times 6}{136}
  \;\approx\; 2.029.
\]
And $2.029 \times B_0 \approx 2.029 \times 25.13 \approx 51.0$ --- not 46.3.
\end{insightbox}

\begin{resultbox}
\textbf{Partial result (promising):}
\[
  \frac{\mu(\Gamma_0(137))}{3}
  \;=\; \frac{138}{3}
  \;=\; 46.
\]
This is within $0.65\%$ of $B_{\mathrm{phenom}} \approx 46.3$.

\textbf{Conjecture A2:}
\[
  B_{\mathrm{phenom}} \;=\; \frac{p+1}{3}
    \bigg|_{p=137}
  \;=\; \frac{138}{3} \;=\; 46,
\]
with the remaining $0.3/46 \approx 0.65\%$ correction from two-loop QED
(already known: the two-loop correction accounts for $\alpha^{-1} = 137 \to
137.036$, an $0.026\%$ effect; a $0.65\%$ correction is too large for
purely two-loop QED).
\end{resultbox}

%──────────────────────────────────────────────────────────────────────────────
\section{The $\mu/3$ Formula: Geometric Origin}
\label{sec:mu3_origin}
%──────────────────────────────────────────────────────────────────────────────

The formula $B_{\mathrm{phenom}} = (p+1)/3$ at $p = 137$ has a natural
geometric interpretation.  The index $\mu = p+1 = 138$ of $\Gamma_0(p)$
in $\mathrm{SL}(2,\mathbb{Z})$ counts the number of inequivalent cusps and
elliptic points weighted appropriately.  The factor $1/3$ appears in the
area formula as $\mathrm{vol}(\mathbb{H}/\mathrm{SL}(2,\mathbb{Z})) = \pi/3$.

In UBT, the ψ-circle partition function sums over $\mathrm{SL}(2,\mathbb{Z})$
orbits of the modular parameter $\tau = it + \psi$.  The natural normalisation
of the kinetic action on the fundamental domain gives a factor of $1/3$ from
the hyperbolic volume element.

\begin{conjecture}[A2: $B_{\mathrm{phenom}}$ from Modular Index]
\label{conj:A2}
The effective loop coefficient is
\begin{equation}
  B_{\mathrm{phenom}} \;=\; \frac{\mu(\Gamma_0(p^*))}{3}
    \;=\; \frac{p^*+1}{3},
\end{equation}
where $p^* = 137$ is the prime attractor selected by UBT stability.
This gives $B_{\mathrm{phenom}} = 46$ exactly, with sub-percent corrections
from subleading geometric terms.
\end{conjecture}

\begin{gapbox}
\textbf{Status of Conjecture~A2:}
\begin{itemize}
  \item Numerical accuracy: $(46.0 - 46.3)/46.3 \approx -0.65\%$.  Close.
  \item Geometric motivation: the modular index $\mu = p+1$ enters the
        area formula for the $\Gamma_0(p)$ quotient --- natural in UBT.
  \item Missing: derivation of the $1/3$ normalisation from $S[\Theta]$
        directly, not from the external area formula.
  \item Missing: explanation of the residual $0.3$ (is it a known
        two-loop correction in disguise?).
\end{itemize}
This is the most numerically promising result of Approach~A2.
It does not close Gap~A but identifies $(p+1)/3$ as the correct
arithmetic structure.
\end{gapbox}

%──────────────────────────────────────────────────────────────────────────────
\section{Conclusion}
\label{sec:conclusion}
%──────────────────────────────────────────────────────────────────────────────

\begin{center}
\begin{tabular}{lll}
  \toprule
  Test & Value & Matches $46.3$? \\
  \midrule
  $g(137)^{3/2}$ & $11^{3/2} \approx 36.5$ & No \\
  $N_{\mathrm{eff}}^{3/2}$ & $12^{3/2} \approx 41.6$ & No ($-10\%$) \\
  $g(137) \times \pi$ & $\approx 34.6$ & No \\
  $\mathrm{vol}(X_0(137))/\pi = \mu/3$ & $46.0$ & \textbf{Yes, to $0.65\%$} \\
  $\mu/3 + \nu_2/4 = 46 + 0.5$ & $46.5$ & \textbf{Yes, to $0.43\%$} \\
  \bottomrule
\end{tabular}
\end{center}

The best result is $\mu(\Gamma_0(137))/3 + \nu_2/4 = 46.5$, within $0.43\%$
of $B_{\mathrm{phenom}} = 46.3$.  Here $\nu_2 = 2$ is the number of elliptic
points of order 2 in $X_0(137)$ (computed by the Kronecker symbol $(-4|137) = 1$).
The full genus formula suggests the complete expression:
\begin{equation}
  B_{\mathrm{phenom}} \;\stackrel{?}{=}\;
  \frac{\mu(\Gamma_0(p^*))}{3} + \frac{\nu_2}{4} + \frac{\nu_3}{3}
  \;=\; 46 + \frac{2}{4} + 0 \;=\; 46.5
  \quad\text{(at $p^* = 137$)},
\end{equation}
off by $0.2$ from $46.3$ --- still $0.43\%$.

\begin{gapbox}
\textbf{Strongest candidate in Approach A2:}
$B_{\mathrm{phenom}} \approx \mu(\Gamma_0(137))/3 + \nu_2(137)/4 = 46.5$
to $0.43\%$ accuracy.  Gap~A is not closed, but the modular-curve geometry
(index and elliptic-point count of $\Gamma_0(137)$) provides the best
numerical candidate found in any approach.
\end{gapbox}

\ifdefined\INCLUDEMODE
\else
  \end{document}
\fi
